\section{Related Work}\label{sec:related}

\paragraph*{Classical spatial indexing}
The R-tree~\cite{guttman1984rtree} and its variants are the canonical
disk-resident spatial indices. For static, read-dominated workloads,
bulk-loading dominates query performance: the Sort-Tile-Recursive (STR)
algorithm~\cite{leutenegger1997str} achieves near-optimal MBR overlap and
underpins PostGIS and Boost.Geometry. M-CTX adopts STR as a strong,
exact, dependency-free baseline.

\paragraph*{Learned spatial indexing}
Following the learned-index paradigm~\cite{kraska2018case}, two-dimensional
methods either project onto a space-filling curve and learn a 1-D
CDF~\cite{wang2019learnedz}, partition space with per-cell models
(LISA~\cite{li2020lisa}), recursively partition with per-leaf models
(RSMI~\cite{qi2020rsmi}), or learn the data layout and partitioning
directly (Flood~\cite{nathan2020flood}, LMSFC~\cite{gao2022lmsfc}). A
recurring but under-examined issue for \emph{range} queries is the
\textbf{centroid-versus-MBR-overlap recall gap}: indexing centroids on a
curve silently drops features whose MBR overlaps the query but whose
centroid does not. Existing one-curve methods close this gap only by an
\emph{ad hoc} global expansion, with no correctness proof and a loose
candidate set. BR-LZ (Section~\ref{sec:brlz}) is, to our knowledge, the
first one-curve learned spatial index to close it with a \emph{proof} of
recall completeness and a provably tighter segment-local expansion.

\paragraph*{Moving-object indexing}
The TPR-tree~\cite{saltenis2000tpr} indexes time-parameterised MBRs;
the B$^{x}$-tree~\cite{jensen2004bx} encodes time-projected positions on a
space-filling curve and indexes them with a plain B$^{+}$-tree, avoiding
costly node updates. Its phase structure matches the AIS $20$\,s cadence
and supports $O(\log N)$ incremental inserts---essential for a feed
ingesting $\sim\!10^{6}$ updates per minute---making it the natural fit
for the neighbour stage~\cite{moon2001hilbert}.

\paragraph*{Distance transforms and geospatial systems}
The Felzenszwalb--Huttenlocher transform~\cite{felzenszwalb2004distance}
computes the exact Euclidean distance transform in linear time via
separable parabolic-envelope passes~\cite{virtanen2020scipy}; the
reference pipeline instead ships a quadratic kernel. General geospatial
engines (PostGIS, GeoMesa) target out-of-core, multi-tenant serving;
M-CTX targets the orthogonal operating point of a single-process,
in-memory accelerator embedded in the training loop, where amortised
per-anchor latency at corpus scale is what matters.

\paragraph*{Context-aware maritime prediction}
EnvShip-Bench~\cite{envshipbench2026} crystallised the environmental and
social context representations we accelerate; earlier work treated AIS as
a trajectory-only signal~\cite{nguyen2018traisformer,zhao2021bilstm} over
raw windows without aligned context~\cite{dma2025pilot}. M-CTX is
orthogonal: it accelerates the context pipeline these methods depend on,
leaving the models and the benchmark unchanged.

\noindent The distinction we draw is sharper than ``learned vs.\
classical.'' Classical structures (R-tree, STR) are exact by
construction; existing learned spatial indices trade that exactness for a
model and recover high recall only empirically, by tuning to a fixed data
distribution. For a retrieval workload feeding a model---where a silently
dropped feature corrupts the training signal---this is the wrong trade. We
therefore ask whether a learned index can retain the build-time and
footprint advantages of the learned paradigm while restoring the
\emph{provable} exactness of the classical one. BR-LZ answers
affirmatively, and Section~\ref{ssec:extent} shows it does so with a
tighter candidate set than the only alternative (a global expansion) that
makes prior learned indices recall-complete.
